
% ============================================================================
\clearpage
\section*{v4.10.0-r19 Source-Only: Status Notes on the Commutant of One Generation, and on Prior Literature for the Doublet and the Octant}
\addcontentsline{toc}{section}{v4.10.0-r19 Source-Only: Status Notes on the Commutant of One Generation, and on Prior Literature for the Doublet and the Octant}
\label{sec:v4100-r19-commutant-literature}
% ============================================================================

\paragraph{Purpose.}
This layer records an exploration of 15 November 2026 on the open item ``a charged carrier'' of v4.10.0-r16.7, that is whether
the generation index of v4.10.0-r15.5 can sit on the frame doublet, and a literature check made for a planned short paper.
The exploration computes all operators on one generation of \(\mathbb C\otimes\HH\otimes\OO\) that commute with the gauge and
Lorentz generators.  Its outcome was fixed by Schur's lemma before the computation; it is not blind and has no
pre-registration.  The literature items attach references to earlier layers and add no new content.  The notes are
\texttt{notes/\allowbreak research/\allowbreak v4\_10\_0\_r19/} and the folders cited there.  Earlier layers are not rewritten.

\noindent\textbf{Source-only status.}
This v4.10.0-r19 layer is source-only.  It does not build a PDF, does not assign a new DOI,
and does not authorize a public release.  The v5.1.0 Version DOI
\href{https://doi.org/10.5281/zenodo.24237112}{10.5281/zenodo.24237112} and the Concept DOI
\href{https://doi.org/10.5281/zenodo.19210728}{10.5281/zenodo.19210728} are unchanged by it.

\subsection*{v4.10.0-r19 status notes}
\label{subsec:v4100-r19-status}

\paragraph{Theoremlet v4.10.0-r19.1 (the commutant of one generation is commutative).}
On \(\mathbb C\otimes\HH\otimes\OO\cong\mathbb C^{32}\) take the generators of the short-paper construction: Lorentz rotations
\(L_q\) and boosts \(iL_q\) on \(\HH\), the weak \(T_a'=C_{\rm ch}T_a\) with \(T_a=-\tfrac i2R_{e_a}\) and \(C_{\rm ch}=\tfrac12(1+iL_{e_7})\),
the hypercharge \(Y/2=Q-\tfrac12+(1-C_{\rm ch})T_3\), and colour as the eight derivations of \(\OO\) that annihilate \(e_7\) (of the
fourteen derivations).  The spectrum of \(Y\) gives the six Weyl multiplets \(Q,L,u^c,d^c,e^c,\nu^c\) with multiplicities
\(12,4,6,6,2,2\).
\begin{enumerate}[leftmargin=2em]
\item[(a)] The complex-linear operators commuting with all generators form a commutative algebra of dimension \(6\), one phase per
multiplet; no antilinear operator commutes with all generators.  The real commutant is \(\mathbb C^6\).
\item[(b)] Hence no pair of anticommuting operators is compatible with the symmetries, and no qubit with three mutually unbiased
bases exists on one generation, whether built from the frame maps or from any other operators.
\item[(c)] Of the frame maps of v4.9.0-r51, \(R=i\cdot1\) lies in the commutant but is central; \(T\) (complex conjugation) and
\(J_B=iT\) do not commute with the Lorentz generators, \(SU(2)'\) and \(Y\), since they map to the conjugate representation.
\item[(d)] With \(n\) copies of the generation the commutant contains \(M_n(\mathbb C)\) acting on the copy index (checked for
\(n=2\) with \(1\otimes\sigma_k\)), so a qubit exists for \(n\ge2\).  On the copy index \(R\) is again a scalar, and \(T\) still does not
commute with the symmetries.
\end{enumerate}
The statement is Schur's lemma for a representation whose irreducible constituents have multiplicity one and are of complex
type; the computation confirms it on the concrete carrier and extends v4.10.0-r16.1 from samples on \(\mathbb C\otimes\OO\) to all
symmetry-compatible operators on one generation.  On the real doublet \((E,B)\) the frame maps confirm v4.10.0-r14.3(d): \(R^2=-1\),
\(T\), \(J_B\), \(R\) span \(\mathfrak{sl}(2,\mathbb R)\), and at most two reflections anticommute; the third unbiased basis needs a
second imaginary unit, which the complex amplitudes of v4.10.0-r16.2 and r16.3 supply on \(\chi\).

\paragraph{Remark v4.10.0-r19.2 (consequence for the charged carrier of r16.7).}
In the Gram reading of v4.10.0-r16.4 the three states are readout states in a doublet \(W\), and the square roots of the masses
are the eigenvalues of their Gram matrix on the generation space \(\mathbb C^3\).  A mass matrix that commutes with the charges
requires the map \(\mathbb C^3\to W\) to act as the identity on the charge part; \(\mathbb C^3\) and \(W\) are then multiplicity
spaces of the gauge representation, and by r19.1 the carrier must contain the generation several times.  The doublet
\(\chi\) of r16.3 can therefore enter a charged carrier only through copies that are present beforehand (the sedenion route of
v4.10.0-r13 has two, under CT-003), and the frame maps of r51 would have to be defined anew on the copy index.  The
reading ``three generations are the three frame maps'' of r15.5 and r16.4 then takes the number of generations as input.
It is not derived on the inventory.  The computations of r15.5(a)--(c), r16.3 and r16.7 on the field doublet stand.

\paragraph{Remark v4.10.0-r19.3 (prior literature for the field doublet).}
For time-harmonic fields the doublet \((E,H)\) has three bases, electric--magnetic, parallel--antiparallel (\(E\pm H\)) and
right--left (helicity), with Stokes-type parameters on a sphere on which duality rotates about the helicity axis and time
reversal acts as a mirror; boosts change the total and the helicity parameter and leave the other two fixed (Golat, Vernon,
Rodr\'iguez-Fortu\~no, arXiv:2405.15718).  In the language of the Companion these are the eigenbases of \(T\), of \(J_B\) (the
compensator ports) and of \(R\).  This answers the open third axis of v4.9.0-r04 from the literature and is the setting of
v4.10.0-r16.2.  The Riemann--Silberstein reading of v4.8.0-r06 (Reading~B) and r07 is textbook material (Bialynicki-Birula and
Bialynicka-Birula, J.~Phys.~A 46 (2013) 053001, arXiv:1211.2655).

\paragraph{Remark v4.10.0-r19.4 (prior literature for the octant and the neutrino phase).}
Three spin-\(\tfrac12\) states from the three mutually unbiased bases, a product of their projectors with phase \(e^{i\pi/4}\), amplitudes
as roots of \(z^3=(1+i)/4\), the charged-lepton form \(1+\sqrt2\cos(2g\pi/3+2/9)\) and the neutrino phase \(2/9+\pi/12\) appear in
C.~A.~Brannen, Found.\ Phys.\ 40 (2010) 1681, arXiv:1006.3114.  The equivalence of \(K=\tfrac23\) with an angle \(\pi/4\) between the
vector of square-root masses and \((1,1,1)\) is due to R.~Foot, arXiv:hep-ph/9402242; the value \(\delta\approx2/9\) is discussed by
P.~\.Zenczykowski, Phys.\ Rev.\ D 86 (2012) 117303 and 87 (2013) 077302.  These references belong to v4.10.0-r12.2, r12.4(c),
r12.5(c), r14.5, r15.4 and r16.7, and to H99, whose pre-registration cites Brannen only through a secondary source.  The
octant triple, its phase \(\pi/4\) and the shift \(\pi/12\) are therefore not results of the Companion; what is specific here is their
placement on the frame maps of the field doublet (r15.5, r16.7), restricted by r19.2.  The two arXiv texts were read through
automated summaries of the full text; formulas and equation numbers are to be confirmed in the PDFs before citation.

\subsection*{v4.10.0-r19 status in the three frames}
\label{subsec:v4100-r19-frames}

\begin{itemize}[leftmargin=2em]
\item M, T: \(R\) acts as a scalar on every complex-linear carrier, so the change \(M\leftrightarrow T\) does not separate states of one
generation; \(T\) maps a generation to its conjugate (r19.1(c)).
\item P: not evaluated.
\end{itemize}

\subsection*{v4.10.0-r19 register entries}
\label{subsec:v4100-r19-register}

\begin{center}
\small
\begin{tabular}{@{}>{\raggedright\arraybackslash}p{0.34\textwidth}>{\raggedright\arraybackslash}p{0.58\textwidth}@{}}
\toprule
statement & status after v4.10.0-r19\\
\midrule
commutant of one generation on \(\mathbb C\otimes\HH\otimes\OO\) & \(\mathbb C^6\), commutative, no antilinear part (r19.1)\\
generation index on the frame doublet (r15.5) & excluded on one generation; possible only with copies set beforehand (r19.2)\\
charged carrier of r16.7 & needs multiplicity at least two; frame maps not defined on the copy index (r19.2)\\
three bases of the field doublet, helicity axis & in the literature (Golat et al.\ 2024) (r19.3)\\
octant triple, phase \(\pi/4\), neutrino shift \(\pi/12\) & in the literature (Brannen 2010; Foot 1994) (r19.4)\\
H99 & unchanged (compatible, not blind); attribution to Brannen confirmed at the primary source\\
H100 / H102 & pending / pending\\
CT-001 to CT-004 & unchanged\\
\bottomrule
\end{tabular}
\end{center}

\subsection*{v4.10.0-r19 ledger and guard update}
\label{subsec:v4100-r19-ledger-guard}

\[
R_{\rm chk}^{v4100\text{-}r19}=(1,1,0)
\]
for (status notes recorded; the commutant computed in the cited script; a carrier with several generations and frame maps
on it not derived).

The layer spends no guard:
\[
R_{\rm glob}^{(v4100\text{-}r19,+)}=(0,1,1,1,1),
\qquad
\Gamma_{\rm guard}^{v4100\text{-}r19}=(1,1,1,1,1),
\qquad
\beta_4=1/256\ \text{unspent}.
\]
OP~1 and OP~3 remain open.  The layer does not imply exact finite \(C5\) zero closure, exact all-readout terminality,
\(C6/T2\), a global photon-clock/RH theorem, or global all-readout closure.
